\section{Evaluating Extensibility (D3)}\label{sec:extensibility}

To evaluate the extensibility requirement (D3), we extend \method to text, a modality it did not previously support, and measure what that extension changed in the existing library.
We then run an intended and an unintended interaction on the new modality with a LoRA-adapted \llama victim: ONION against a textual backdoor attack, and DP-SGD against the same attack.

\subsection{Provisions for Extensibility}

\method organizes modules by risk, so a new attack, defense, or metric is implemented under the namespace of the risk it targets, and the modules already present in that risk are unchanged (Figure~\ref{fig:overview}).
Each risk defines an abstract base class that fixes the entry point its defenses expose: defenses against~\ref{poison} implement a robust-training method, and defenses against~\ref{meminf} implement a private-training method.
A conformance test asserts that every defense in the package exposes the entry point defined for its risk, and fails the build otherwise.
New datasets enter through the \texttt{AmuletDataset} class, and every helper function accepts configurable parameters, so \method integrates into an existing PyTorch pipeline.

\method provides a contribution guide and uses standard, modern Python tooling: \texttt{uv}\footnote{\url{https://docs.astral.sh/uv/}} for development environments, \texttt{ruff}\footnote{\url{https://docs.astral.sh/ruff/}} and \texttt{pyright}\footnote{\url{https://microsoft.github.io/pyright/}} for code quality checks, and \texttt{pre-commit}\footnote{\url{https://pre-commit.com/}} hooks that run those checks before each commit.
Figures~\ref{fig:extending-risk} and~\ref{fig:extending-model} in the appendix show the steps for adding a new risk and a new model architecture.

\subsection{Extending \method to Text}

Supporting text required four new modules.
\texttt{TextTensorDataset} represents a tokenized text dataset.
\texttt{HFCausalLM} wraps a pretrained causal language model with low-rank adapters~\cite{hu2021LoRALowRankAdaptation} and a linear classification head.
\texttt{TextBadNets} implements the textual form of the BadNets attack~\cite{gu2019BadNetsEvaluatingBackdooring}, inserting a trigger word into a training record and relabeling that record to the target label.
\texttt{ONION} implements a perplexity-based input purification defense against~\ref{poison}~\cite{qi2021ONIONSimpleEffective}.
The four modules total approximately $1{,}000$ lines.
Each implements the abstract base class for its risk, so the metrics already implemented for~\ref{poison} consume their outputs without modification.
The second study required no new module at all: \method already implements DP-SGD~\cite{abadi2016DeepLearningDifferential} as a defense against~\ref{meminf}, and it applies to the LoRA-adapted victim without modification.

Beyond the package export lists, the extension modified one existing file.
That modification added the text dataset class, and widened the \texttt{AmuletDataset} modality field to admit text alongside images and tabular records.
No existing risk, attack, defense, or metric changed.
Supporting an entirely new modality therefore cost four new modules and one widened type, and left every module already in the library intact.
A privacy defense written and tested against vision and tabular victims now measures a poisoning attack on a three-billion-parameter LLM, and neither the defense nor the metrics that consume its output were touched to make that possible.

This follows from how we represent a text dataset.
We define \texttt{TextTensorDataset} as a subclass of PyTorch's \texttt{TensorDataset} over token identifiers and labels, so existing consumers receive the type they already expect.
Poisoning attacks keep their return type, and training loops that iterate over data and label pairs run without modification.

The conformance test shaped ONION's interface.
ONION purifies inputs, and the interface it would naturally expose differs from the retraining interface that the other poisoning defenses share.
The test requires the shared entry point, so ONION retrains the victim on purified data through it, and a user substitutes one poisoning defense for another without changing the surrounding pipeline.

\subsection{Experimental Setup}

We use \sstwo~\cite{socher2013RecursiveDeepModels} for both studies and measure the effect of ONION and DP-SGD on~\ref{poison}.
\sstwo provides $67{,}349$ training records, and we use its $872$-record validation split as the test set, since the labels of the official test split are withheld.
The victim is \llama~\cite{llamateammetaai2024Llama3Herd}, adapted with LoRA at rank $8$ and a linear classification head over the pooled final-token representation, trained in full precision.
\texttt{TextBadNets} inserts the trigger \texttt{cf} at a random position in a poisoned record and assigns that record the target label.
Random placement follows the protocol of the original evaluation and prevents ONION from identifying the trigger by position alone.

Each study compares three victims.
$\modelstd$ trains on the clean dataset and establishes the baseline.
$\modelpois$ trains on the poisoned dataset with no defense.
$\modeldef$ trains on the same poisoned dataset with the defense applied.
We report test accuracy ($Acc_{te}$) on the clean test set, and attack success rate ($ASR$), the fraction of triggered test inputs that the victim assigns to the target class.
ONION purifies both the training data and the inputs at inference, so we evaluate $\modeldef$ on purified inputs.
DP-SGD applies during training only, so we evaluate its victim on unmodified inputs.
We sweep the poison rate from $0.01\%$ to $5\%$, corresponding to $6$ and $3{,}367$ records.
All experiments are repeated five times with mean and standard error reported.

\begin{table*}[htb]
    \setlength\tabcolsep{4pt}
    \centering
    \footnotesize
    \caption{Poisoning~(\ref{poison}) versus two defenses on \sstwo with a
    LoRA-adapted \llama victim. Each block adds one defense to the poisoned model
    $\modelpois$; both share the clean baseline $\modelstd$. ONION drives $ASR$ down
    at every poison rate, while DP-SGD suppresses the backdoor only at small poison
    counts. $^{*}$This cell collapsed to the majority class and does not measure the
    defense.}%
    \label{tab:textbadnets_interactions}
    \begin{tabular}{ c c c c c c c }
        \bottomrule

        \toprule
        \multirow{2}{*}{\textbf{Poisoned}} & \multirow{2}{*}{\textbf{\# Poisoned}} & \multirow{2}{*}{\textbf{Defense}} & \multicolumn{2}{c}{\textbf{Undefended} ($\modelpois$)} & \multicolumn{2}{c}{\textbf{Defended} ($\modeldef$)} \\
        & & & $Acc_{te}$ & $ASR$ & $Acc_{te}$ & $ASR$ \\
        \midrule
        Baseline ($\modelstd$) & $-$ & $-$ & 95.76~$\pm$~0.17 & $-$ & $-$ & $-$ \\
        \midrule
        \multicolumn{7}{c}{\textbf{ONION} (intended interaction)} \\
        \midrule
        0.01\% & 6 & \multirow{5}{*}{ONION} & 95.39~$\pm$~0.27 & 19.21~$\pm$~5.68 & 93.51~$\pm$~0.17 & 7.85~$\pm$~0.81 \\
        0.1\% & 67 &  & 95.48~$\pm$~0.37 & 98.27~$\pm$~1.16 & 94.22~$\pm$~0.17 & 8.04~$\pm$~0.72 \\
        1\% & 673 &  & 95.60~$\pm$~0.30 & 99.95~$\pm$~0.05 & 93.83~$\pm$~0.31 & 8.08~$\pm$~1.06 \\
        2\% & 1346 &  & 95.57~$\pm$~0.17 & 99.95~$\pm$~0.05 & 93.94~$\pm$~0.48 & 6.78~$\pm$~0.60 \\
        5\% & 3367 &  & 95.00~$\pm$~0.53 & 100.00~$\pm$~0.00 & 93.17~$\pm$~0.39 & 8.32~$\pm$~1.50 \\
        \midrule
        \multicolumn{7}{c}{\textbf{DP-SGD} (unintended interaction)} \\
        \midrule
        \multirow{2}{*}{0.1\%} & \multirow{2}{*}{67} & $\epsilon = 1$ & \multirow{2}{*}{94.04} & \multirow{2}{*}{99.77} & 71.67 & 29.44 \\
         &  & $\epsilon = 8$ &  &  & 70.41 & 33.64 \\
        \multirow{2}{*}{1\%} & \multirow{2}{*}{673} & $\epsilon = 1$ & \multirow{2}{*}{95.53} & \multirow{2}{*}{99.77} & 68.81 & 43.46 \\
         &  & $\epsilon = 8$ &  &  & 72.13 & 49.07 \\
        \multirow{2}{*}{5\%} & \multirow{2}{*}{3367} & $\epsilon = 1$ & \multirow{2}{*}{93.23} & \multirow{2}{*}{100.00} & 50.92$^{*}$ & 100.00$^{*}$ \\
         &  & $\epsilon = 8$ &  &  & 68.92 & 97.20 \\
        \bottomrule

        \toprule
    \end{tabular}
\end{table*}


\subsection{Intended Interaction:~\ref{poison} and ONION}

Results for both studies are shown in Table~\ref{tab:textbadnets_interactions}.
The attack reaches $98.3\%$ $ASR$ at $67$ poisoned records and $100\%$ above that, while $Acc_{te}$ stays within a point of the $95.76\%$ baseline.
ONION reduces $ASR$ to $6.8$--$8.3\%$ at a cost of roughly two points of $Acc_{te}$, and removed the trigger from every triggered input in every run.
The original evaluation of ONION on \sstwo reports near-total attack success without a defense, and single to low double digit $ASR$ under the defense at a small accuracy cost~\cite{qi2021ONIONSimpleEffective}.
Our results on \llama are consistent with these values.

The sweep also shows two effects that the original evaluation does not report.
First, $ASR$ under ONION holds near $8\%$ across a $560$-fold change in the poison count.
Once ONION removes the trigger, the poison count no longer affects the outcome.
Second, the attack has a minimum scale.
At $67$ poisoned records out of $67{,}349$ it succeeds nearly every time, while at $6$ records mean $ASR$ falls to $19.2\%$ with a standard error of $5.7$, and per-seed values span $2.1\%$ to $33.9\%$.
Backdooring an LLM of this size therefore takes very few poisoned records, and the attack is unreliable below that point.

\subsection{Unintended Interaction:~\ref{poison} and DP-SGD}

DP-SGD~(\ref{dpsgd}) protects against~\ref{meminf} and is not designed to protect against~\ref{poison}.
It nonetheless suppresses the backdoor while the poison count is small.
At $67$ poisoned records it reduces $ASR$ from $99.8\%$ to $29.4$--$33.6\%$, and at $673$ records it reduces $ASR$ to $43.5$--$49.1\%$.
The effect then disappears: at $3{,}367$ poisoned records the privately trained victim reaches $97.2\%$ $ASR$ at $\epsilon=8$, leaving the attack as effective as it is against an undefended victim.

This pattern follows the group-privacy bound on poisoning a private learner.
An $\epsilon$-private learner forces an adversary to modify at least $\frac{1}{\epsilon}\log\tau$ records to reduce the attack cost by a factor of $\tau$, so the protection holds for small poison counts and decays as that count grows~\cite{ma2019DataPoisoningDifferentiallyPrivate}.
DP-SGD bounds the influence that any single training record exerts on the trained model, and its protection for a group of records weakens as that group grows.
A backdoor planted with $67$ records is a small set of outliers that clipping and noise suppress.
A backdoor planted with $3{,}367$ records forms a subpopulation whose gradient direction many records share, and a bound on individual influence does not constrain it.

Private training costs $24$--$27$ points of $Acc_{te}$ in every configuration we ran.
At $67$ poisoned records this buys a $70$-point reduction in $ASR$, and at $3{,}367$ records it buys none.
Measuring this interaction at a single poison rate would mischaracterize the defense in either direction.

The bound certifies protection for small poison counts, but empirical studies of DP-SGD reach conclusions that depend on the setting.
Measured against a non-private baseline, private models can be more susceptible than undefended ones, because the accuracy that privacy removes and the accuracy that poisoning removes fall on the same records~\cite{jagielski2021DoesDifferentialPrivacy}.
A separate analysis attributes the poisoning benefit of DP-SGD to gradient clipping alone, with the useful configurations sitting at privacy budgets too loose to certify~\cite{hong2020EffectivenessMitigatingData}.
These studies evaluate convex models or small networks, and none reports a backdoor on an LLM.
Our result adds a point they do not cover: on a three-billion-parameter victim, DP-SGD suppresses the backdoor at $\epsilon=1$ when the poison count is small, a meaningful privacy budget, though at a large utility cost.
The sign and size of this interaction depend on the model, the attack, and the privacy budget, and cannot be read off from any single study.

One configuration, $\epsilon=1$ at $3{,}367$ poisoned records, collapsed to predicting the target class for every input, which yields a degenerate $ASR$ of $100\%$ alongside chance accuracy.
We mark that cell in Table~\ref{tab:textbadnets_interactions} and exclude it from the analysis above.

\begin{takeaway}
Extending \method to text and LLMs cost four new modules and reused a privacy defense unmodified.
That defense suppressed the backdoor only at small poison counts.
\end{takeaway}
